\section{Summary and Future Work}

This paper develops a Euclidean, \emph{Principia}-style treatment of the Kepler problem using finite-step constructions, tangent/triangle geometry, and Newton's ultimate-ratio viewpoint. For the inverse problem, the proof chain across Sections~2--5 is explicit and modular. Section~2 establishes the dynamical bridge: from constant areal speed, the acceleration direction is central, and two local lemmas convert a geometric ratio into the inverse-square force form. Section~3 introduces the auxiliary-circle/affine transport layer; the tangent-transfer and matching-drop construction reduces local deflection to a conic-invariant ratio, culminating in Inverse Problem Proof~1 for the ellipse (\Cref{subsec:proof1-ellipse-affine}). A more universal proof architecture then generalizes to all conic types: conic-specific geometry computes the local ratio, while the Section~2 conversion step remains universal.

Within that inverse-problem architecture, Proof~2 and Proof~3 play different roles. Proof~2 is based on direct displacement computation, identifying the constant sagitta-to-chord-square ratio. Proof~3 uses a more direct hodograph route: first establish that conic motion gives a circular hodograph, then use \(\Delta \mathbf v/\Delta t\) to recover acceleration, which is inverse-square in radius.

For the forward problem, we label the two forward routes as Proof~1F and Proof~2F, with variant closures denoted by Proof~2F'. In Proof~1F we keep the hodograph-circle theorem in view but proceed in a discrete Newtonian style. The infinitesimal \(\Delta t\)-scaled hodograph circle is translated in parallel as it glides along the orbit: its shape and orientation are preserved while its center shifts from step to step. This moving-circle picture provides a direct geometric bridge from velocity-space circularity to conic recovery in configuration space. 

Our second forward route, Proof~2F, is fully geometric. From the initial condition and fixed physical constants, the hodograph is determined; after rotation and scaling, the remaining steps are ruler-and-compass constructions. The method is pointwise and unique: from local data (the velocity \(\mathbf v\), with both magnitude and direction, together with invariant \(L\)), we set the arm length \(d:=FH=L/v\), identify \(H\), draw \(t_A\perp FH\), and impose the centripetal-direction condition \(AF\parallel OV\). These constraints determine the orbital point \(A\) uniquely in each local configuration. Repeating the construction yields the full orbit. A key invariant is a fixed second focus \(F'\) for ellipse/hyperbola, or a fixed directrix line \(\mathcal D\) for parabola; this certifies the conic class. Variants denoted by Proof~2F' are mainly technical Euclidean line-construction variants that demonstrate flexibility of the method, but are not required for the main architecture.

We also include multiple geometric realizations of the hodograph (and rotated/scaled proxies), including different choices of velocity origin. Several parabola constructions appear to be less emphasized in the existing literature. The force-center auxiliary-circle viewpoint used throughout is also less common than directrix-circle-centered expositions. The infinitesimal \(\Delta t\)-scaled moving-hodograph-circle picture appears to be a useful and potentially original pedagogical viewpoint. More broadly, we intentionally provide a denser geometric presentation than the minimalist style of the \emph{Principia}, supported by modern construction software and figures, to make relation chains explicit and accessible.

Finally, the main limitation remains geometric uniformity across conic types: Inverse Problem Proof~1 (the ellipse affine-circle route) is especially clean, but we have not yet found an equivalent proof of the same type for hyperbola and parabola, which currently require additional case-specific constructions.

\subsection*{Acknowledgment}
This study was motivated by appreciation of geometric beauty and by sustained work on geometric approaches and their extension to broader generalizations, especially through ellipse properties, auxiliary-circle structures, and their related identities. A major precursor was the author's work on an elementary geometry problem for middle-school students, documented in \cite{cryptodogares2026geometrydegreefreedom}.
That direction led to a deeper reading of Newton's \emph{Principia} treatment of space-time Euclidean geometry, and then to possible new exploration paths for both inverse and forward problems.

The author is deeply grateful to the long geometric tradition represented by Euclid, Galileo, Kepler, Apollonius, Newton, Leibniz, Cauchy, Hamilton, Maxwell, Feynman, and Bernoulli, and to modern contributors including Goodstein \& Goodstein, Chandrasekhar, Markowsky, Derbes, van Haandel--Heckman, and Cari\~nena--Ra\~nada--Santander. The manuscript is intended as a pedagogical account for geometric tradition and its application to mechanics.

\subsection*{Code and Tooling Disclosure}
Code, \LaTeX\ sources, and geometric construction files are available at:
\url{https://github.com/CryptoDogAres/AlternativeKeplerToNewton/}

For long-term citation stability, the best practice is to include both the repository URL and the arXiv identifier.

Figures were produced with GeoGebra Classic 6. Writing and editing assistance used OpenAI ChatGPT, Codex, and Prism (including skill-assisted proofreading and \LaTeX\ editing support). Iterative drafting and paper editing were performed in Visual Studio Code.
